\documentclass[11pt,a4paper]{article}
\usepackage[utf8]{inputenc}
\usepackage[T1]{fontenc}
\usepackage[french]{babel}
\usepackage{geometry}
\usepackage{amsmath}
\usepackage{amssymb}
\usepackage{enumitem}
\usepackage{titlesec}
\usepackage{xcolor}

\geometry{margin=2.5cm}

% Couleurs personnalisées
\definecolor{primarycolor}{RGB}{31, 78, 121}
\definecolor{secondarycolor}{RGB}{220, 100, 50}

% Configuration des titres
\titleformat{\section}{\large\bfseries\color{primarycolor}}{}{0pt}{}[\titlerule]
\titleformat{\subsection}{\normalsize\bfseries\color{secondarycolor}}{}{0pt}{}

\setlist[itemize]{label=\textbullet, leftmargin=*, nosep, after=\vspace{0.5em}}

\title{\textbf{\Large Progression de mathématiques --- Classes de 4\textsuperscript{ème} (2026-27)}}
\author{}
\date{}

\begin{document}

\maketitle
\vspace{-3em}

\section{Trimestre 1}

\subsection{Séquence 1 : Addition et Soustraction de Nombres relatifs}
\textbf{Objectifs d’apprentissage :}
\begin{itemize}
    \item Maîtriser les règles opératoires d’addition et de soustraction sur les nombres décimaux relatifs.
    \item Découvrir et appliquer les règles de multiplication et de division des nombres relatifs (règle des signes).
    \item Utiliser la notation simplifiée en supprimant les signes d’addition et les parenthèses superflues.
    \item Effectuer des enchaînements de calculs et des comparaisons pour traiter et résoudre des problèmes concrets.
\end{itemize}

\textbf{Prolongements possibles :}
\begin{itemize}
    \item Approche historique : l'introduction des quantités négatives à travers les notions de dettes et de fortunes chez le mathématicien indien Brahmagupta (VII\textsuperscript{e} siècle) et la résistance de l'Occident chrétien qui les qualifiait encore de "nombres absurdes" au XVI\textsuperscript{e} siècle.
\end{itemize}

\textbf{Automatismes :}
\begin{itemize}
    \item Déterminer instantanément le signe d'un produit comportant plusieurs facteurs négatifs.
    \item Effectuer des calculs mentaux rapides d'additions et de soustractions de relatifs (ex : $-7 - 3 = -10$ ; $-5 + 8 = 3$).
\end{itemize}

\subsection{Séquence 2 : Théorème de Pythagore (partie 1)}
\textbf{Objectifs d’apprentissage :}
\begin{itemize}
    \item Connaître et utiliser l'égalité de Pythagore dans un triangle rectangle.
    \item Mobiliser la liste des carrés parfaits de 1 à 144 et comprendre la notion de racine carrée.
    \item Déterminer la longueur d'un côté d’un triangle rectangle à partir de la longueur des deux autres.
    \item Mener des raisonnements géométriques simples et s’initier à la rédaction d’une démonstration de calcul de longueur.
\end{itemize}

\textbf{Prolongements possibles :}
\begin{itemize}
    \item Étude des triplets pythagoriciens sur les tablettes d'argile babyloniennes (Plimpton 322, vers -1800 avant J.-C.) et découverte de la corde à 13 nœuds utilisée par les arpenteurs égyptiens (les "tendeurs de corde") pour tracer des angles droits sur les chantiers.
\end{itemize}

\textbf{Automatismes :}
\begin{itemize}
    \item Donner instantanément le carré d’un nombre inférieur ou égal à 12, ou la racine carrée d'un carré parfait usuel.
    \item Isoler mathématiquement le terme manquant dans une égalité du type $a^2 = b^2 + c^2$.
\end{itemize}

\subsection{Séquence 3 : Arithmétique}
\textbf{Objectifs d’apprentissage :}
\begin{itemize}
    \item Consolider la notion de nombre premier et déterminer les nombres premiers inférieurs ou égaux à 100.
    \item Décomposer un nombre entier en un produit de facteurs premiers, à la main ou à l’aide d’un logiciel.
    \item Trouver des diviseurs communs à deux entiers et appréhender la notion de PGCD de manière intuitive.
    \item Simplifier des fractions de manière experte pour les rendre irréductibles.
    \item Modéliser et résoudre des problèmes algorithmiques ou de la vie courante mettant en jeu la divisibilité (engrenages, cycles mécaniques).
\end{itemize}

\textbf{Prolongements possibles :}
\begin{itemize}
    \item Introduction aux nombres parfaits (égaux à la somme de leurs diviseurs propres, comme 6 et 28) et étude historique du mécanisme d'Anticythère, la plus vieille machine à calculer astronomique basée sur des engrenages de nombres premiers.
\end{itemize}

\textbf{Automatismes :}
\begin{itemize}
    \item Appliquer instantanément les critères de divisibilité par 2, 3, 4, 5, 9 et 10.
    \item Reconnaître si un petit nombre entier est premier ou non.
\end{itemize}

\subsection{Séquence 4 : Théorème de Pythagore (partie 2) Réciproque}
\textbf{Objectifs d’apprentissage :}
\begin{itemize}
    \item Connaître et formuler la réciproque et la contraposée du théorème de Pythagore.
    \item Utiliser l'égalité de Pythagore pour démontrer qu’un triangle est rectangle (perpendicularité) ou ne l’est pas.
    \item Structurer une démonstration géométrique rigoureuse en séparant le calcul du côté le plus long de la somme des carrés des deux autres côtés.
\end{itemize}

\textbf{Prolongements possibles :}
\begin{itemize}
    \item Analyse du puzzle de Bhaskara (mathématicien indien du XII\textsuperscript{e} siècle) qui démontre le théorème de Pythagore visuellement par un simple découpage d'aires géométriques sous le seul mot d'ordre : "Regarde !".
\end{itemize}

\textbf{Automatismes :}
\begin{itemize}
    \item Identifier sans erreur l'hypoténuse potentielle d'un triangle à partir de la donnée de ses trois longueurs.
    \item Rédiger les phrases clés d'une démonstration ("D'une part...", "D'autre part...").
\end{itemize}

\subsection{Séquence 5 : Rationnels / Fractions Addition soustraction partie 1}
\textbf{Objectifs d’apprentissage :}
\begin{itemize}
    \item Approfondir la somme et la différence de deux nombres rationnels à partir d'exemples de partages concrets.
    \item Calculer l’addition ou la soustraction de deux fractions ayant le même dénominateur.
    \item Effectuer des calculs de sommes et de différences lorsque le dénominateur de l'une est un multiple du dénominateur de l'autre.
\end{itemize}

\textbf{Prolongements possibles :}
\begin{itemize}
    \item Interdisciplinarité Musique et Maths : l'étude de la gamme pythagoricienne où les intervalles musicaux (octave, quinte, quarte) sont construits à partir de rapports fractionnaires de longueurs de cordes vibrantes (1/2, 2/3, 3/4).
\end{itemize}

\textbf{Automatismes :}
\begin{itemize}
    \item Trouver rapidement un dénominateur commun simple entre deux fractions (recherche de multiples communs).
    \item Simplifier le résultat d'un calcul fractionnaire.
\end{itemize}

\subsection{Séquence 6 : Triangles égaux}
\textbf{Objectifs d’apprentissage :}
\begin{itemize}
    \item Définir deux triangles égaux comme deux triangles superposables.
    \item Connaître et utiliser les propriétés caractéristiques associées (égalité des longueurs des côtés homologues et des mesures des angles homologues).
    \item Faire le lien entre les cas d’égalité des triangles et la construction géométrique d’un triangle à partir de longueurs et/ou de mesures d’angles données.
\end{itemize}

\textbf{Prolongements possibles :}
\begin{itemize}
    \item Étude historique des techniques de triangulation maritime au XVII\textsuperscript{e} siècle, où l'égalité et la similitude des triangles permettaient aux navigateurs de calculer la distance d'un navire par rapport à la côte sans traverser l'eau.
\end{itemize}

\textbf{Automatismes :}
\begin{itemize}
    \item Identifier et associer les sommets, côtés et angles homologues de deux triangles superposables.
\end{itemize}

\subsection{Séquence 7 : Rationnels / Fractions Multiplication partie 2}
\textbf{Objectifs d’apprentissage :}
\begin{itemize}
    \item Définir et appliquer la règle de multiplication de deux nombres rationnels (produit des numérateurs et produit des dénominateurs).
    \item Développer des stratégies de simplification de fractions en décomposant les facteurs avant d'effectuer le produit.
    \item Résoudre des problèmes de la vie courante impliquant la prise d'une fraction d'une autre quantité fractionnaire.
\end{itemize}

\textbf{Prolongements possibles :}
\begin{itemize}
    \item Les fractions égyptiennes : étude du Papyrus Rhind (-1650 avant J.-C.) montrant comment les scribes écrivaient toutes les fractions uniquement sous forme de sommes de fractions unitaires (de numérateur 1) pour distribuer équitablement les rations de pain.
\end{itemize}

\textbf{Automatismes :}
\begin{itemize}
    \item Multiplier un entier par une fraction.
    \item Repérer des facteurs identiques au numérateur et au dénominateur pour simplifier un produit avant calcul (ex : $3/5 \times 5/7 = 3/7$).
\end{itemize}

\section{Trimestre 2}

\subsection{Séquence 8 : Puissances}
\textbf{Objectifs d’apprentissage :}
\begin{itemize}
    \item Définir la puissance d’un nombre pour des exposants entiers relatifs (positifs et négatifs).
    \item Connaître et manipuler les puissances de 10 ainsi que les règles opératoires de calcul associées.
    \item Maîtriser la notation scientifique des nombres décimaux.
    \item Passer d’une représentation d’un nombre à une autre et effectuer des calculs numériques simples impliquant des puissances.
\end{itemize}

\textbf{Prolongements possibles :}
\begin{itemize}
    \item Les ordres de grandeur en astronomie : calcul de la distance Terre-Soleil (unité astronomique) ou de la vitesse de la lumière en notation scientifique, et initiation à la loi de Titius-Bode sur les distances des planètes.
\end{itemize}

\textbf{Automatismes :}
\begin{itemize}
    \item Exprimer des grands ou petits nombres sous forme de puissances de 10 (ex : $1~000~000 = 10^6$ ; $0,001 = 10^{-3}$).
    \item Convertir un nombre en notation scientifique.
\end{itemize}

\subsection{Séquence 9 : Rationnels / Fractions Division partie 3}
\textbf{Objectifs d’apprentissage :}
\begin{itemize}
    \item Définir la notion d'inverses d'un nombre rationnel non nul et distinguer l'inverse de l'opposé.
    \item Établir et appliquer la règle de division de deux nombres rationnels : diviser par une fraction revient à multiplier par son inverse.
    \item Mener des calculs numériques combinant les quatre opérations sur les fractions en respectant les priorités opératoires.
\end{itemize}

\textbf{Prolongements possibles :}
\begin{itemize}
    \item Introduction aux paradoxes de Zénon d'Élée (le paradoxe d'Achille et de la tortue) pour observer de manière philosophique comment la division infinie du temps et de l'espace se traduit par une suite de fractions décroissantes.
\end{itemize}

\textbf{Automatismes :}
\begin{itemize}
    \item Donner immédiatement l'inverse d'un nombre entier ou d'une fraction.
\end{itemize}

\end{document}
